\input{estilo}
% Propuesta, opción 1: la cola interna y los notifiers. Lo que un notifier no
% pudo entregar vuelve a la misma cola (el lazo rojo), y los mensajes cuyos
% suscriptores responden (verdes) quedan detrás de los pendientes (rojos).
% Misma entrada y mismas posiciones que la opción 2 (cola-de-retry), para
% leerlas como par. Agregado del texto: los workers, el lazo y los pies.
\begin{document}
\begin{figura}\begin{tikzpicture}
  \newcommand\msj[2]{\node[draw=#1, fill=#1!14, minimum size=4.6mm, inner sep=0pt] at #2 {};}
  \node[cola, minimum height=40mm, minimum width=11mm] (q) at (2.6,0) {};
  \foreach \c [count=\i] in {acento, segundo, acento, segundo, acento, segundo}
    {\msj{\c}{({1.3+0.55*(\i-1)},0)}}
  \draw[msg] (-0.4,0) -- (q.west);
  \foreach \d in {0.3, 0.15} \node[caja, fill=white, minimum width=22mm, minimum height=14mm] at (7.2+\d,-\d) {};
  \node[caja, minimum width=22mm, minimum height=14mm] (w) at (7.2,0) {notifiers};
  \draw[msg] (q.east) -- (w.west);
  \draw[msg, draw=acento] (w.north) -- ++(0,0.65) -| node[pos=0.25, above, etiq, text=acento] {lo que no pudo entregar vuelve a la misma cola} (0.2,0.1);
  \node[etiq, text=suave] at (3.9,-1.15) {los que funcionan esperan detrás de los pendientes};
  \node[anchor=west] at (-0.4,-1.9) {\textbf{opción 1:} volver a encolar};
  % la leyenda de colores, igual en las dos opciones
  \node[anchor=west, font=\footnotesize] at (-0.4,-2.5) {%
    \tikz[baseline=-0.6ex]\node[draw=segundo, fill=segundo!14, minimum size=3.5mm] {};\ su suscriptor responde \quad
    \tikz[baseline=-0.6ex]\node[draw=acento, fill=acento!14, minimum size=3.5mm] {};\ pendiente};
\end{tikzpicture}\end{figura}
\end{document}
